
\subsection{Leggere visivamente una matrice RowMask23}
\label{sec:guida-visiva-rowmask}

L'equazione~\eqref{eq:rowmask23} descrive un singolo coefficiente, ma la struttura risulta più intuitiva osservando contemporaneamente tutte le sue componenti. La Figura~\ref{fig:tre-componenti-rowmask} usa un ritaglio con le righe e colonne $0,1,8,9$ della stessa matrice $128\times128$: il ritaglio serve a vedere gli indici, non rappresenta un blocco autonomo con parametri diversi.

\begin{figure}[htbp]
\centering
\begin{tikzpicture}[x=8.5mm, y=-6.5mm]
\begin{scope}[xshift=0mm]
\node[text=steraGrigio] at (0.5,-.45) {0}; \node[text=steraGrigio] at (-.45,0.5) {0};
\node[text=steraGrigio] at (1.5,-.45) {1}; \node[text=steraGrigio] at (-.45,1.5) {1};
\node[text=steraGrigio] at (2.5,-.45) {8}; \node[text=steraGrigio] at (-.45,2.5) {8};
\node[text=steraGrigio] at (3.5,-.45) {9}; \node[text=steraGrigio] at (-.45,3.5) {9};
\node[font=\small\bfseries] at (2,-1.4) {$D$: diagonale};
\draw[steraGrigio!60, fill=steraBlu!20] (0,0) rectangle ++(1,1); \node[font=\footnotesize] at (0.5,0.5) {$d_{0}$};
\draw[steraGrigio!60, fill=white] (1,0) rectangle ++(1,1); \node[font=\footnotesize,text=steraGrigio] at (1.5,0.5) {0};
\draw[steraGrigio!60, fill=white] (2,0) rectangle ++(1,1); \node[font=\footnotesize,text=steraGrigio] at (2.5,0.5) {0};
\draw[steraGrigio!60, fill=white] (3,0) rectangle ++(1,1); \node[font=\footnotesize,text=steraGrigio] at (3.5,0.5) {0};
\draw[steraGrigio!60, fill=white] (0,1) rectangle ++(1,1); \node[font=\footnotesize,text=steraGrigio] at (0.5,1.5) {0};
\draw[steraGrigio!60, fill=steraBlu!20] (1,1) rectangle ++(1,1); \node[font=\footnotesize] at (1.5,1.5) {$d_{1}$};
\draw[steraGrigio!60, fill=white] (2,1) rectangle ++(1,1); \node[font=\footnotesize,text=steraGrigio] at (2.5,1.5) {0};
\draw[steraGrigio!60, fill=white] (3,1) rectangle ++(1,1); \node[font=\footnotesize,text=steraGrigio] at (3.5,1.5) {0};
\draw[steraGrigio!60, fill=white] (0,2) rectangle ++(1,1); \node[font=\footnotesize,text=steraGrigio] at (0.5,2.5) {0};
\draw[steraGrigio!60, fill=white] (1,2) rectangle ++(1,1); \node[font=\footnotesize,text=steraGrigio] at (1.5,2.5) {0};
\draw[steraGrigio!60, fill=steraBlu!20] (2,2) rectangle ++(1,1); \node[font=\footnotesize] at (2.5,2.5) {$d_{8}$};
\draw[steraGrigio!60, fill=white] (3,2) rectangle ++(1,1); \node[font=\footnotesize,text=steraGrigio] at (3.5,2.5) {0};
\draw[steraGrigio!60, fill=white] (0,3) rectangle ++(1,1); \node[font=\footnotesize,text=steraGrigio] at (0.5,3.5) {0};
\draw[steraGrigio!60, fill=white] (1,3) rectangle ++(1,1); \node[font=\footnotesize,text=steraGrigio] at (1.5,3.5) {0};
\draw[steraGrigio!60, fill=white] (2,3) rectangle ++(1,1); \node[font=\footnotesize,text=steraGrigio] at (2.5,3.5) {0};
\draw[steraGrigio!60, fill=steraBlu!20] (3,3) rectangle ++(1,1); \node[font=\footnotesize] at (3.5,3.5) {$d_{9}$};
\end{scope}
\node[font=\Large] at (4.75,2) {$+$};
\begin{scope}[xshift=47mm]
\node[text=steraGrigio] at (0.5,-.45) {0}; \node[text=steraGrigio] at (-.45,0.5) {0};
\node[text=steraGrigio] at (1.5,-.45) {1}; \node[text=steraGrigio] at (-.45,1.5) {1};
\node[text=steraGrigio] at (2.5,-.45) {8}; \node[text=steraGrigio] at (-.45,2.5) {8};
\node[text=steraGrigio] at (3.5,-.45) {9}; \node[text=steraGrigio] at (-.45,3.5) {9};
\node[font=\small\bfseries] at (2,-1.4) {$B$: base XOR};
\draw[steraGrigio!60, fill=white] (0,0) rectangle ++(1,1); \node[font=\footnotesize,text=steraGrigio] at (0.5,0.5) {0};
\draw[steraGrigio!60, fill=white] (1,0) rectangle ++(1,1); \node[font=\footnotesize] at (1.5,0.5) {$b_{1}$};
\draw[steraGrigio!60, fill=steraVerde!25] (2,0) rectangle ++(1,1); \node[font=\footnotesize] at (2.5,0.5) {$b_{8}$};
\draw[steraGrigio!60, fill=white] (3,0) rectangle ++(1,1); \node[font=\footnotesize] at (3.5,0.5) {$b_{9}$};
\draw[steraGrigio!60, fill=white] (0,1) rectangle ++(1,1); \node[font=\footnotesize] at (0.5,1.5) {$b_{1}$};
\draw[steraGrigio!60, fill=white] (1,1) rectangle ++(1,1); \node[font=\footnotesize,text=steraGrigio] at (1.5,1.5) {0};
\draw[steraGrigio!60, fill=white] (2,1) rectangle ++(1,1); \node[font=\footnotesize] at (2.5,1.5) {$b_{9}$};
\draw[steraGrigio!60, fill=steraVerde!25] (3,1) rectangle ++(1,1); \node[font=\footnotesize] at (3.5,1.5) {$b_{8}$};
\draw[steraGrigio!60, fill=steraVerde!25] (0,2) rectangle ++(1,1); \node[font=\footnotesize] at (0.5,2.5) {$b_{8}$};
\draw[steraGrigio!60, fill=white] (1,2) rectangle ++(1,1); \node[font=\footnotesize] at (1.5,2.5) {$b_{9}$};
\draw[steraGrigio!60, fill=white] (2,2) rectangle ++(1,1); \node[font=\footnotesize,text=steraGrigio] at (2.5,2.5) {0};
\draw[steraGrigio!60, fill=white] (3,2) rectangle ++(1,1); \node[font=\footnotesize] at (3.5,2.5) {$b_{1}$};
\draw[steraGrigio!60, fill=white] (0,3) rectangle ++(1,1); \node[font=\footnotesize] at (0.5,3.5) {$b_{9}$};
\draw[steraGrigio!60, fill=steraVerde!25] (1,3) rectangle ++(1,1); \node[font=\footnotesize] at (1.5,3.5) {$b_{8}$};
\draw[steraGrigio!60, fill=white] (2,3) rectangle ++(1,1); \node[font=\footnotesize] at (2.5,3.5) {$b_{1}$};
\draw[steraGrigio!60, fill=white] (3,3) rectangle ++(1,1); \node[font=\footnotesize,text=steraGrigio] at (3.5,3.5) {0};
\end{scope}
\node[font=\Large] at (10.3,2) {$+$};
\begin{scope}[xshift=94mm]
\node[text=steraGrigio] at (0.5,-.45) {0}; \node[text=steraGrigio] at (-.45,0.5) {0};
\node[text=steraGrigio] at (1.5,-.45) {1}; \node[text=steraGrigio] at (-.45,1.5) {1};
\node[text=steraGrigio] at (2.5,-.45) {8}; \node[text=steraGrigio] at (-.45,2.5) {8};
\node[text=steraGrigio] at (3.5,-.45) {9}; \node[text=steraGrigio] at (-.45,3.5) {9};
\node[font=\small\bfseries] at (2,-1.4) {$R$: correzioni};
\draw[steraGrigio!60, fill=white] (0,0) rectangle ++(1,1); \node[font=\footnotesize,text=steraGrigio] at (0.5,0.5) {0};
\draw[steraGrigio!60, fill=white] (1,0) rectangle ++(1,1); \node[font=\footnotesize,text=steraGrigio] at (1.5,0.5) {0};
\draw[steraGrigio!60, fill=white] (2,0) rectangle ++(1,1); \node[font=\footnotesize,text=steraGrigio] at (2.5,0.5) {0};
\draw[steraGrigio!60, fill=white] (3,0) rectangle ++(1,1); \node[font=\footnotesize,text=steraGrigio] at (3.5,0.5) {0};
\draw[steraGrigio!60, fill=steraArancio!25] (0,1) rectangle ++(1,1); \node[font=\footnotesize] at (0.5,1.5) {$e_{1,19}$};
\draw[steraGrigio!60, fill=white] (1,1) rectangle ++(1,1); \node[font=\footnotesize,text=steraGrigio] at (1.5,1.5) {0};
\draw[steraGrigio!60, fill=white] (2,1) rectangle ++(1,1); \node[font=\footnotesize,text=steraGrigio] at (2.5,1.5) {0};
\draw[steraGrigio!60, fill=steraArancio!25] (3,1) rectangle ++(1,1); \node[font=\footnotesize] at (3.5,1.5) {$e_{1,0}$};
\draw[steraGrigio!60, fill=steraArancio!25] (0,2) rectangle ++(1,1); \node[font=\footnotesize] at (0.5,2.5) {$e_{8,6}$};
\draw[steraGrigio!60, fill=white] (1,2) rectangle ++(1,1); \node[font=\footnotesize,text=steraGrigio] at (1.5,2.5) {0};
\draw[steraGrigio!60, fill=white] (2,2) rectangle ++(1,1); \node[font=\footnotesize,text=steraGrigio] at (2.5,2.5) {0};
\draw[steraGrigio!60, fill=white] (3,2) rectangle ++(1,1); \node[font=\footnotesize,text=steraGrigio] at (3.5,2.5) {0};
\draw[steraGrigio!60, fill=white] (0,3) rectangle ++(1,1); \node[font=\footnotesize,text=steraGrigio] at (0.5,3.5) {0};
\draw[steraGrigio!60, fill=white] (1,3) rectangle ++(1,1); \node[font=\footnotesize,text=steraGrigio] at (1.5,3.5) {0};
\draw[steraGrigio!60, fill=white] (2,3) rectangle ++(1,1); \node[font=\footnotesize,text=steraGrigio] at (2.5,3.5) {0};
\draw[steraGrigio!60, fill=white] (3,3) rectangle ++(1,1); \node[font=\footnotesize,text=steraGrigio] at (3.5,3.5) {0};
\end{scope}
\end{tikzpicture}

\caption{Rappresentazione delle tre componenti della medesima matrice RowMask23: diagonale $D$, base XOR $B$ e correzioni sparse $R$. I quattro indici mostrati sono un sottoinsieme del blocco reale $128\times128$. La parte evidenziata nella matrice centrale riutilizza una stessa coordinata $b_8$.}
\label{fig:tre-componenti-rowmask}
\end{figure}

Il primo termine, $D$, lascia libera ogni posizione $(i,i)$ e vale zero altrove. Il secondo, $B$, assegna il medesimo valore $b_s$ a tutte le posizioni con $i\oplus j=s$, purché $s\ne0$. Queste posizioni non formano necessariamente una diagonale geometrica nel senso ordinario: per questa ragione si parla di \emph{diagonali generalizzate} rispetto all'operazione XOR. Il terzo termine, $R$, interviene soltanto nelle posizioni scelte dalla funzione $\sigma(i,k)$.

Si osservi in particolare la coordinata $b_8$. Nella posizione $(0,8)$ vale $0\oplus8=8$, quindi il coefficiente riceve $b_8$. Anche nella posizione $(1,9)$ vale $1\oplus9=8$. Tuttavia la riga $1$, per $k=0$, ha offset
\[
\sigma(1,0)=1+((13\cdot0+7\cdot1)\bmod127)=8.
\]
Pertanto $(1,9)$ è una posizione dotata di correzione, mentre la riga $0$ è stata costruita senza correzioni. Si ottiene
\[
W_{0,8}=b_8,\qquad W_{1,9}=b_8+e_{1,0}.
\]
Questo esempio evidenzia l'idea centrale: \emph{condivisione} del parametro di base e \emph{adattamento locale} mediante una coordinata addizionale.

\subsection{Esempio numerico di ricostruzione}
\label{sec:esempio-numerico-rowmask}

Per rendere esplicita la ricostruzione, assegniamo soltanto a scopo illustrativo i valori
\[
d_1=0{,}60,\qquad b_8=0{,}25,\qquad e_{1,0}=-0{,}10.
\]
Non sono pesi estratti dal modello addestrato: sono numeri scelti per consentire la verifica a mano. Le tre caselle della Tabella~\ref{fig:lettura-casella-rowmask} permettono di distinguere il trattamento della diagonale, della sola base e della base con correzione.

\begin{table}[htbp]
\centering
\small
\begin{tabular}{cccl}
\toprule
Casella $(i,j)$ & $s=i\oplus j$ & Origine & Peso risultante \\
\midrule
$(1,1)$ & $0$ & diagonale & $d_1 = 0{,}60$ \\
$(0,8)$ & $8$ & base, riga senza correzioni & $b_8 = 0{,}25$ \\
$(1,9)$ & $8$ & base e correzione $k=0$ & $b_8 + e_{1,0} = 0{,}25 - 0{,}10 = 0{,}15$ \\
\bottomrule
\end{tabular}
\caption{Lettura di tre caselle di un blocco RowMask23, con valori dimostrativi. Le caselle $(0,8)$ e $(1,9)$ condividono la stessa coordinata di base $b_8$; solo la seconda riceve anche una correzione.}
\label{fig:lettura-casella-rowmask}
\end{table}

La ricostruzione fornisce
\[
W_{1,1}=0{,}60,\qquad W_{0,8}=0{,}25,\qquad
W_{1,9}=0{,}25-0{,}10=0{,}15.
\]
Il fatto che $W_{0,8}$ e $W_{1,9}$ utilizzino entrambi $b_8$ non costringe i loro valori a coincidere: la correzione $e_{1,0}$ consente di modificare selettivamente il secondo. Il modello non memorizza i tre coefficienti come tre pesi densi indipendenti; salva $d_1$, $b_8$ ed $e_{1,0}$ nelle rispettive regioni del vettore di coordinate.

\subsection{Perché la maschera non è una potatura della matrice}
\label{sec:non-potatura-rowmask}

La parola \emph{maschera} può suggerire erroneamente che il modello memorizzi una matrice quasi tutta nulla. In RowMask23 non è così: la maggior parte degli elementi fuori diagonale dipende da una coordinata della base XOR e può assumere un valore diverso da zero. La parte sparsa riguarda \emph{solo le correzioni}. Si possono distinguere tre quantità:
\begin{itemize}
\item \num{16384} coefficienti della matrice densa $128\times128$, che è la rappresentazione espansa;
\item \num{2921} correzioni della matrice $R$, collocate in $23$ posizioni per ciascuna delle $127$ righe non nulle della maschera;
\item \num{3176} coordinate complessive del blocco, sommando $128$ valori diagonali, $127$ valori della base XOR e le \num{2921} correzioni.
\end{itemize}

Le \num{3176} coordinate rappresentano il $19{,}38\%$ degli elementi densi: la riduzione dei gradi di libertà del solo blocco è dunque circa $80{,}62\%$. Questo non è, però, il risparmio complessivo del Transformer STERA 30M. Il modello completo aggiunge scale, guadagni, fattori a basso rango, vettori di posizione, normalizzazioni e bias, che vengono conteggiati separatamente nel capitolo sull'architettura. Inoltre, \emph{rango pieno} e \emph{arbitrarietà dei coefficienti} sono due proprietà differenti: la famiglia RowMask23 può avere rango pieno e restare al tempo stesso un sottoinsieme proprio delle matrici dense.

\subsection{Dalla ricostruzione al gradiente}
\label{sec:esempio-gradiente-rowmask}

Lo stesso esempio chiarisce come si addestra una rappresentazione condivisa. Sia $G=\partial\mathcal{L}/\partial W$ il gradiente rispetto alla matrice materializzata in una singola operazione. Se una certa coordinata compare in più caselle, la regola della catena \emph{somma} i contributi delle caselle corrispondenti.

Per la coordinata $b_8$ sono interessate tutte le posizioni $(i,i\oplus8)$, non soltanto le due mostrate nel ritaglio:
\[
\frac{\partial\mathcal{L}}{\partial b_8}
=\sum_{i=0}^{127}G_{i,i\oplus8}.
\]
Viceversa, per $d_1$ ed $e_{1,0}$ il gradiente proviene da una sola posizione:
\[
\frac{\partial\mathcal{L}}{\partial d_1}=G_{1,1},
\qquad
\frac{\partial\mathcal{L}}{\partial e_{1,0}}=G_{1,9}.
\]
Questa distinzione spiega perché un'ottimizzazione standard sulle coordinate è possibile: dopo aver differenziato la perdita rispetto ai valori utilizzati nel prodotto, il trainer accumula le derivate nelle posizioni del vettore compatto e aggiorna queste ultime. La matrice ricostruita non acquisisce gradi di libertà indipendenti durante l'ottimizzazione.

\subsection{Significato della ricostruzione esatta}
\label{sec:esattezza-ricostruzione}

La denominazione inglese STERA include \emph{Reconstruction}. La proprietà che il nome indica è precisa. Una matrice $128\times128$ che ha la struttura RowMask23, cioè che appartiene all'immagine di $\Phi$, si codifica con le sue \num{3176} coordinate e si ricostruisce da esse in modo esatto: per la linearità e l'iniettività della Sezione~\ref{sec:linearita-iniettivita}, la corrispondenza fra le coordinate e le matrici della famiglia è biunivoca, e la ricostruzione non richiede di stimare coefficienti mancanti. La proprietà riguarda le matrici della famiglia: una matrice densa qualunque non ammette in generale una rappresentazione RowMask23, perché la famiglia ha dimensione \num{3176}, inferiore a \num{16384}.

Questa corrispondenza è ciò che serve all'addestramento. L'ottimizzatore non modifica una matrice densa da proiettare poi sulla famiglia: modifica le coordinate, e a ogni vettore di coordinate corrisponde una e una sola matrice. Ne seguono tre conseguenze. L'addestramento può raggiungere ogni matrice della famiglia, quindi conserva per intero l'espressività che la struttura RowMask23 mette a disposizione. Non esistono coordinate ridondanti, cioè direzioni di aggiornamento che lascino invariata la matrice, quindi ogni coordinata contribuisce al modello. Infine il gradiente rispetto alle coordinate (Sezione~\ref{sec:gradiente-coordinate}) è calcolato attraverso la stessa $\Phi$ usata nel passo in avanti, quindi ciò che l'ottimizzazione migliora è esattamente la matrice $\Phi(c)$ che il modello usa.

La stessa esattezza vale fra addestramento e inferenza: la ricostruzione eseguita dai kernel di inferenza coincide bit per bit con quella dell'addestramento (Sezione~\ref{sec:kernel-rowmask}). Il modello che si esegue è quindi lo stesso che si è addestrato, senza approssimazioni nel passaggio dall'una all'altra fase.
